\section{Conclusion and Future Work}
\label{sec:conclusion}

In this paper, we conducted a rigorous, scientifically honest empirical study and post-mortem of on-device model compression for merged multi-task networks. Specifically, we investigated \textbf{ZipMerge}, a framework designed to co-optimize layer-wise merging coefficients and magnitude-pruning boundaries at test-time using an unsupervised minimum entropy objective on tiny calibration sets. 

Evaluating a compact Vision Transformer (\texttt{vit\_tiny\_patch16\_224}) backbone across four highly disparate visual classification tasks (MNIST, FashionMNIST, CIFAR-10, and SVHN), our experiments exposed critical empirical boundaries that have extensive implications for physical system deployment. We demonstrated that all merged models suffer from catastrophic representational collapse, performing at the level of random guessing (10\% to 14\% absolute accuracy) due to severe representational and spatial weight conflicts. We analyzed why the simple, unoptimized decoupled baseline, Prune-then-Merge (P-then-M), consistently outperformed joint test-time optimization under high task conflict, revealing that pre-merging pruning acts as a crucial spatial regularizer that removes orthogonal parameter noise. Furthermore, we detailed how unconstrained unsupervised test-time adaptation on tiny calibration sets triggers the \emph{Overfitting-Optimizer Paradox}, where entropy minimization drives down training loss while destroying generalizable features.

Rather than trying to spin a failure as a success, our work establishes a highly valuable and realistic baseline study of model merging and pruning limits in the wild. We translated these empirical failures into actionable architectural guidelines, advocating for:
\begin{enumerate}[leftmargin=*]
    \item Restricting merging to domain-aligned task groups where representations remain compatible.
    \item Leveraging low-rank PEFT adapters (such as LoRA) fine-tuned on a massive pre-trained foundation model to preserve base representations.
    \item Employing explicit test-time regularizers to prevent the Overfitting-Optimizer Paradox from destroying network generalizability.
\end{enumerate}

Future research will focus on investigating these regularized adapter-merging schemes to enable truly robust and efficient multi-task acceleration on edge hardware. Specifically, we will explore advanced PEFT subspace methods such as \textbf{Orthogonal LoRA Initialization} (to pre-allocate non-interfering update pathways across tasks) and \textbf{Orthogonal Procrustes Alignment} (which we formulate and analyze in Section~\ref{sec:peft_study}, demonstrating its lightweight computational complexity on commodity hardware). Fusing aligned adapters minimizes representation mismatch in the low-rank manifold, actively closing the performance gap between merged parameter-efficient adapters and specialized dense experts. 

We envision a promising co-design of ZipMerge with post-training quantization (PTQ) to achieve multiplicative compression gains. Under an INT8 or INT4 weight quantization scheme, the non-differentiable step of quantization can be naturally integrated into the ZipMerge loop alongside the pruning mask. Specifically, during test-time adaptation, we define a uniform quantization step size $q_{\text{step}}$ based on the maximum absolute weight of the merged parameters:
\begin{equation}
    q_{\text{step}} = \frac{2 \max(|\theta_{\text{merged}}(\Lambda)|)}{2^b - 1}
\end{equation}
where $b$ represents the bit-width (e.g., $b=4$ or $b=8$). The joint non-differentiable composition of dynamic magnitude pruning and uniform quantization is then formulated as:
\begin{equation}
    \theta_{\text{joint}}(\Lambda) = \text{round}\left(\frac{\theta_{\text{sparse}}(\Lambda)}{q_{\text{step}}}\right) \times q_{\text{step}}
\end{equation}
where $\theta_{\text{sparse}}(\Lambda) = M \odot \theta_{\text{merged}}(\Lambda)$ represents the sparse parameters under the dynamic pruning mask $M$. Under first-order optimization, the gradient of the unsupervised adaptation loss $\mathcal{L}_{\text{TTA}}$ with respect to the continuous coefficients $\Lambda$ is passed directly through both the rounding and pruning operators via the Identity-pass Straight-Through Estimator (STE):
\begin{equation}
    \frac{\partial \mathcal{L}_{\text{TTA}}}{\partial \Lambda} \approx \frac{\partial \mathcal{L}_{\text{TTA}}}{\partial \theta_{\text{joint}}} \cdot \frac{\partial \theta_{\text{merged}}}{\partial \Lambda}
\end{equation}
This mathematical integration illustrates how ZipMerge's dual optimization engines (STE and ES) are uniquely suited to optimize $\Lambda$ through joint pruning and quantization boundaries without requiring expensive gradient backpropagation or full fine-tuning datasets. This co-design would allow practitioners to compress multi-task merged models to extreme sub-structural footprints (e.g., 80\% sparse, 4-bit quantized) for direct deployment on commodity edge NPUs. We explicitly note that this post-training quantization co-design is currently proposed as a purely theoretical framework; its empirical validation and the associated hardware execution latency studies on physical edge NPUs are reserved as exciting directions for future work.

Additionally, while the current ZipMerge framework relies on dynamic magnitude pruning applied abruptly at the start of the adaptation, future work will explore alternative pruning schedules to improve optimization stability. Instead of masking to the target sparsity $p$ at the very first step, a progressive or scheduled magnitude pruning approach can be employed. Under this paradigm, the target sparsity $p_t$ at adaptation step $t$ follows a smooth schedule from $p_0 = 0.0$ to the final target $p$ over the total $T$ adaptation steps. We compare several possible scheduling formulations and discuss their physical impact on the test-time adaptation loop:
\begin{enumerate}[leftmargin=*]
    \item \textbf{Linear Schedule ($p^{\text{lin}}_t = p \cdot \frac{t}{T}$):} Prunes a constant fraction of parameters at each step. This provides a highly predictable and steady trajectory, but may cause premature representation loss in early steps before the merging coefficients $\Lambda$ have settled into a collaborative joint basin.
    \item \textbf{Cubic Schedule ($p^{\text{cub}}_t = p \cdot \left[1 - \left(1 - \frac{t}{T}\right)^3\right]$):} Prunes extremely slowly in the early stages, keeping the network mostly dense so that the coefficients can rapidly converge to a joint flat region, and then aggressively prunes the remaining parameters in the final steps. This allows the network to establish strong multi-task pathways before enforcing high sparsity.
    \item \textbf{Sine/Cosine Schedule ($p^{\text{sin}}_t = p \cdot \sin\left(\frac{\pi t}{2 T}\right)$):} Follows a smooth concave profile, providing a balanced trade-off where the pruning rate is moderate in the beginning and tapers off at the end, ensuring maximum stability near convergence.
\end{enumerate}
These scheduled approaches allow the merging coefficients $\Lambda$ to adapt smoothly to a gradually sparsified parameter space, preventing optimization shocks, avoiding transductive gradient divergence, and preserving crucial base representations.

Furthermore, future work will explore integrating alternative, modern zero-shot pruning criteria directly into the ZipMerge co-optimization loop to guide dynamic mask generation. Specifically, instead of relying purely on weight magnitudes, we can co-optimize the coefficients $\Lambda$ under \textbf{activation-weighted magnitude metrics} inspired by Wanda \cite{sun2023wanda} or \textbf{second-order curvature-based metrics} derived from SparseGPT \cite{frantar2023sparsegpt}. For a merged weight matrix $W_{\text{merged}}(\Lambda)$ at layer $l$, the Wanda-style pruning importance score for the $j$-th column and $i$-th row is defined as:
\begin{equation}
    I_{i, j}(\Lambda) = |[W_{\text{merged}}(\Lambda)]_{i, j}| \cdot \|X_j\|_2
\end{equation}
where $\|X_j\|_2$ is the L2-norm of the incoming activation channel across the calibration batch. By scaling weight magnitudes with active calibration-set activation norms, Wanda-style pruning evaluates the functional importance of individual parameter paths under the current adapted coefficients $\Lambda$. Under this joint paradigm, as the blending coefficients $\Lambda$ are updated, the activation distributions across layers shift dynamically, causing the pruning importance scores $I_{i, j}(\Lambda)$ and the resulting sparse mask to adapt co-operatively. In high-conflict regimes, this activation-weighted magnitude pruning can actively route and preserve vital, overlapping task pathways that standard magnitude pruning might otherwise destructively zero out, representing an exciting and promising next step for research.

Finally, while our empirical studies focus on Vision Transformers and CNNs, model merging has become a key technique for combining large autoregressive language models (LLMs) such as LLaMA and Mistral. Extending the ZipMerge co-optimization framework to autoregressive sequence generation tasks represents an exciting and highly impactful future direction. In this context, rather than computing entropy over image classification probability vectors, test-time adaptation can be driven by minimizing the Shannon entropy of the vocabulary-level next-token prediction distributions (or by optimizing next-token perplexity) over a small set of generic, unlabeled text prompts. Integrating ZipMerge's dynamic pruning and coefficient search with low-rank adapters (e.g., QLoRA) of LLMs would enable extreme, hardware-friendly on-device compression of massive generative models.

\section*{Reproducibility Statement and Code Availability}
To promote transparency, reproducibility, and research in test-time adaptive model compression, the complete PyTorch-compatible source code of ZipMerge, including configuration schemas, evaluation scripts, fine-tuned expert checkpoints, and calibration image indices, is fully open-sourced and publicly accessible at \url{https://github.com/anonymous/zipmerge} under the MIT License.
